\begin{proposicion}[Order of the parameters and errors]\label{prop:cortes-uniformes}
With the historical band and the parameters below, there is a constant \(C_\delta\) common to all
admissible cut-offs. At scale \(n\) the errors in Appendices \ref{app:signed}, \ref{app:arrival},
\ref{app:interior} and \ref{app:weights} decompose as follows. We write \(o_b(1)\) for a quantity tending
to zero as \(n\to\infty\) for each fixed \(b\).
\begin{enumerate}
\item In the signed identity and the arrival remainder, the replacement of \(g^\flat\) and its
exponential costs at most \(C_\delta b+o_b(1)\). The exits, the substitution of approximate minors, the
activation tail, and the Taylor remainders in the stopped identity are \(o_b(1)\).
\item In the interior mean the following terms are \(o_b(1)\): the summed insertion error, the change of
law for the increments, the substitution of the uniform walk averages, and their Riemann sum over levels.
\item In the insertion of a pair, the Taylor remainders in the action and the prefactor contribute
\(o_b(1)\) after summation over pairs. Replacing their coefficients by the eight state integrals at that
same cut-off also costs \(o_b(1)\).
\end{enumerate}
\end{proposicion}
\begin{proof}
The parameters are chosen in the order
\[
 \delta\ \text{and the class margins}\ \longrightarrow\ a\ \longrightarrow\ \rho\ \longrightarrow\
 \alpha\ \longrightarrow\ \delta_*\ \longrightarrow\ b\ \longrightarrow\ n_0(b),
\]
and each choice depends only on the earlier ones.
Fix the band of~\cite[§5.4]{tc} and its \(\delta\), and reserve the common margins of diffusion,
vacancy, gap, load, and exit \(c_*>2\beta\). The class transition fixes \(C_{\rm pres}\) before the
cut-off. These margins fix the radii of the minors and their fields, the threshold \(\eta\), and the base
\(C_1\) of both the deterministic and integrated counts up to \(8\log n\). The exponential reserve for
long prescriptions is chosen before the Taylor truncation fraction and the threshold \(\eta\). Choose
\(0<a<1\) with \(a\log(\max(2,C_1))<1/4\), then \(\rho\ge2\) with \(a\log\rho>3\), and finally
\(\alpha<\alpha_{\rm geo}\) with \(2C\alpha<\eta/(2\rho^2)\) and with the slack of the outer class and
the rook balancing. Choose \(\delta_*<1\) with \(C_1\delta^{-2}\delta_*^{1/3}<1\) and
\(C_2\delta^{-2}\delta_*<1\) for the two series in~\eqref{eq:llegada-nu}. We may now choose any
\(0<b<\min(\delta/2,\delta_*)\) small enough that \(x_0<1\) and the exponential inequality of
Lemma~\ref{lem:llegada-integrada} holds. The proof of Lemma~\ref{lem:prescripcion-larga} shows where this
inequality absorbs the positive part \(\beta j\) of the weight of a long line. Finally, \(n_0(b)\)
ensures \(j\ge\sqrt n,\ v^2/j\le\delta_*\), \(C(4v+12)\le j^{1/3}\) up to \(v=8\log n\), and all class
bounds for the child. The global radii do not increase after the load has been fixed.

For the first item, both series of types are geometric once \(b/\delta_*\) has been extracted. Their
bases and the bounded multiplier for transfer to \(\mu\) depend only on the choices made before \(b\),
increasing \(n_0(b)\) so that \(e^{|A-R_Q(n)|}\le2\). The remaining cost is
\[
 O_b\!\left(nj^{-8}+n^{-1}+n(1+j)e^{-2j}+ne^{-c_b n}
           +n e^{-c_\delta n+D_\delta n^{2/3}}\right)=o_b(1).
\]
The quadratic remainders in the signed identity are \(O_b(n^{-1})\) at scale \(n\). In the second item,
the cubic residuals cost \(O_b(n\sum_{m>j}m^{-3})\). The centered increments and the change of law cost
\(O_b(n(\sum_{m>j}m^{-4})^{1/2})\). Both tend to zero. If the uniform approximation to the mean has error
\(\epsilon_n(b)\to0\), its sum costs \(n\epsilon_n(b)\sum_{m>j}m^{-2}=o_b(1)\). To obtain this
approximation, first fix the walk length, take \(n\to\infty\), and then use the geometric tail of the
resolvent on \([b,1]\). Continuity on this interval makes the Riemann sum error tend to zero. In the
third item the pointwise error in the insertion ratio is \(O_b(j^{-3})\). Each pair carries mass
\(O_b(j^{-2})\), and there are \(O(j^3)\) pairs. The total cost is \(O_b(n/j^2)=o_b(1)\). Accidental
coincidences in the truncated walks cost \(O_{b,N}(1/n)\), and the longer terms are bounded by
\(C_b\sum_{q>N}(1+q)^2r_b^{2q-O(1)}\), \(r_b<1\). These limits justify the substitutions at fixed cut-off.
\end{proof}
Subsequently, comparing the boundary at cut-off \(b\) with its final value gives
\(\eta(b)=\limsup_n|\mathcal M_n(b)-(P_T+J-S)|\to0\) as \(b\) decreases, as used in
Appendix~\ref{app:closing}. This last error belongs to the limit of the boundary profile in \(b\). The
proposition uses \(o_b(1)\) only for the limit in \(n\).
